%% ma: L11-B09-0017
%% lop: 11
%% chuong: 3
%% bai: 9
%% dang: 11.9.6
%% loai: TN4
%% muc: NB
%% dap_an: "C"
%% nguon: "de-on-tap-giua-ki-1-toan-11-nam-2026-2027-truong-thpt-ngo-quyen-da-nang.pdf (M09, Đề ôn tập giữa kì 1 Toán 11 THPT Ngô Quyền, Đà Nẵng), Mã đề 0114, Phần I Câu 8"
%% kien_thuc: "Bài 8 (mẫu số liệu ghép nhóm), Bài 9 (các số đặc trưng đo xu thế trung tâm)"
%% trang_thai: cho-duyet
%% ly_do: "Dạng mới đề xuất 11.9.6"
\begin{ex}
Thống kê quãng đường đi bộ mỗi ngày (đơn vị: km) của một nhóm người cao tuổi thu được bảng số liệu ghép nhóm sau:
\begin{center}\small\begin{tabular}{|l|c|c|c|c|c|c|}\hline
Quãng đường (km) & $[1;2)$ & $[2;3)$ & $[3;4)$ & $[4;5)$ & $[5;6)$ & $[6;7)$\\\hline
Số người & $6$ & $15$ & $28$ & $12$ & $4$ & $6$\\\hline
\end{tabular}\end{center}
Giá trị đại diện của nhóm $[3;4)$ bằng
\choice{$3$}{$4$}{\True $3{,}5$}{$1$}
\loigiai{Giá trị đại diện của nhóm $[3;4)$ là $\dfrac{3+4}2=3{,}5$.}
\end{ex}
